% Generated draft: requires mathematical and source review.
\documentclass[11pt]{article}
\usepackage[margin=25mm]{geometry}
\usepackage{fontspec}
\setmainfont{Latin Modern Roman}
\usepackage{amsmath,amssymb,mathtools,unicode-math}
\setmathfont{Latin Modern Math}
\usepackage{enumitem,needspace,xurl,hyperref}
\setlength{\parindent}{0pt}
\setlength{\parskip}{5pt}
\setlength{\emergencystretch}{3em}
\newcommand{\N}{\mathbb{N}}
\newcommand{\Z}{\mathbb{Z}}
\newcommand{\Q}{\mathbb{Q}}
\newcommand{\R}{\mathbb{R}}
\newcommand{\C}{\mathbb{C}}
\newcommand{\F}{\mathbb{F}}
\newcommand{\A}{\mathbb{A}}
\newcommand{\PP}{\mathbb{P}}
\newcommand{\E}{\mathbb{E}}
\newcommand{\eps}{\varepsilon}
\newcommand{\dd}{\,\mathrm{d}}
\begin{document}
\section*{UnsolvedMath ID: 30001798}
\textbf{Source problem code:} catalogue entry
\subsection{Definitions and mathematical statement}
Let \(q=p^r\), where \(p\) is a prime and \(r\geq 1\), and let \(\F_q\) be the finite field with \(q\) elements. For a smooth proper scheme \(X\) over \(\F_q\), let \(K_i(X)\) denote its Quillen algebraic \(K\)-group in degree \(i\). An abelian group \(G\) is torsion if every \(g\in G\) is annihilated by some positive integer, which may depend on \(g\).

\textbf{Parshin's conjecture, as formulated in the supplied record.} For every finite field \(\F_q\), every smooth proper scheme \(X\) over \(\F_q\), and every integer \(i>0\),
\[
K_i(X)\otimes_{\Z}\Q=0.
\]
Equivalently,
\[
\forall\alpha\in K_i(X)\quad \exists n\in\Z_{>0}\quad n\alpha=0.
\]
The assertion imposes no vanishing condition on \(K_0(X)\) and does not assert that the positive-degree groups are finite.

The source also describes an equivalent formulation using rational motivic cohomology vanishing outside a predicted degree. It does not specify the cohomological degree, weight, or indexing convention, so that formulation cannot be reproduced precisely from the supplied record. The statement above retains the source's smooth proper hypothesis; the dated literature assessment instead refers to smooth projective varieties, and their relationship is not explained in the supplied evidence.
\subsection{Short English statement}
Parshin's conjecture here asks whether every element of every positive-degree algebraic \(K\)-group of a smooth proper scheme over a finite field has finite order. Equivalently, all these groups vanish after tensoring with \(\Q\). The supplied motivic reformulation and literature assessment require clarification.
\subsection{Sources}
\begin{itemize}
\item Supplied local UnsolvedMath registry record: permanent ID \(30001798\), exact source problem code \texttt{catalogue entry}, title \emph{Parshin's Conjecture}, collection \emph{Oberwolfach Reports Open Problems}. The record names Parshin's conjecture but supplies no individual proposer. Its proposed-year field is \(2011\); this is not independently established as the conjecture's original date.
\item \emph{New Directions in Algebraic \(K\)-Theory} (2011), Oberwolfach Reports, pp.~1469--1509, DOI \texttt{10.4171/owr/2011/27}. Supplied report reference: \url{https://doi.org/10.4171/owr/2011/27}. No precise problem page or passage is supplied.
\item Supplied literature reference labelled ``rational Parshin theorem'': \url{https://arxiv.org/abs/2011.05544}.
\item Supplied literature reference labelled ``surface cases'': \url{https://doi.org/10.1016/j.jpaa.2023.107518}.
\item Literature assessment dated exactly \texttt{2026-08-21}: the supplied curation labels the problem partially solved and states that Parshin vanishing is proved rationally and in important special cases, while the full conjecture for all smooth projective varieties over finite fields remains open. This is supplied dated curation derived from the OWR record, not an independently verified literature assessment. The supplied background expressly requests independent verification of time-sensitive status and cited claims.
\end{itemize}
\end{document}
